\chapter{Multiple antennas}\label{ch:mimo}

% ------------------------------------------------------------------
\section{The flat-fading MIMO channel}
With $N_t$ transmit and $N_r$ receive antennas and flat fading, one channel use is
\begin{equation}
  \vect y=\mat H\vect x+\vect w,\qquad \mat H=\frac{1}{\sqrt{N_t}}\mat G,\quad G_{ij}\sim\CN(0,1),\quad \vect w\sim\CN(\vect0,\sigma^2\mat I).
\end{equation}
The normalization by $\sqrt{N_t}$ keeps the total transmitted power independent of the number of
antennas. Every receive antenna then sees unit average signal power, and the SNR set in the
\blk{awgn} block is the SNR \emph{per receive antenna}. Antenna correlation follows the
Kronecker model with exponential correlation matrices,
\begin{equation}
  \mat H=\frac{1}{\sqrt{N_t}}\,\mat R_r^{1/2}\mat G\,\mat R_t^{1/2},\qquad [\mat R]_{ij}=\rho^{|i-j|}.
\end{equation}

\begin{incode}
Multi-stream signals have shape $(B,N_t,n)$.
\begin{itemize}
  \item \blk{stream_demux} distributes symbols round-robin over the streams and sets
  \meta{spatial_rate}$=N_t$, so that $E_b/N_0$ accounts for the $N_t$ symbols sent per channel
  use (\cref{eq:ebn0}).
  \item \blk{mimo_channel} draws $\mat H$ per frame (\prm{block}) or per symbol (\prm{fast}) and
  publishes it on \code{csi} (\meta{csi_type}\code{="mimo"}).
  \item \blk{mimo_detector} and \blk{alamouti_decoder} use it, and \blk{stream_mux} restores a
  single stream.
\end{itemize}
\end{incode}

% ------------------------------------------------------------------
\section{Detection for spatial multiplexing}
\subsection{Zero forcing}
$\mat W_{\mathrm{ZF}}=(\mat H\Her\mat H)^{-1}\mat H\Her$ separates the streams exactly,
$\hat{\vect x}=\vect x+\mat W\vect w$. Stream $k$ then has noise variance
$\sigma^2[(\mat H\Her\mat H)^{-1}]_{kk}$, which becomes large when $\mat H$ is ill-conditioned. For
$N_t=N_r$ each stream has only one degree of diversity.

\subsection{MMSE and its bias}
The linear MMSE estimator is $\mat W=(\mat H\Her\mat H+\sigma^2\mat I)^{-1}\mat H\Her$ (unit-power
symbols). Its output is biased: $\E[\hat x_k|x_k]=b_kx_k$ with $b_k=(\mat W\mat H)_{kk}<1$.
\begin{proposition}
For the MMSE filter, the residual interference plus noise of stream $k$ has variance
$b_k(1-b_k)$. After unbiasing ($\hat x_k/b_k$), the noise variance is therefore $(1-b_k)/b_k$,
and the post-detection SINR is $b_k/(1-b_k)$.
\end{proposition}
\begin{proof}
For the MMSE filter, the error covariance is
$\mat E=\E[(\hat{\vect x}-\vect x)(\hat{\vect x}-\vect x)\Her]=\mat I-\mat W\mat H$, so the MSE of
stream $k$ is $1-b_k$. Write $\hat x_k=b_kx_k+\nu_k$ with $\nu_k$ uncorrelated with $x_k$ (an
orthogonality property of the MMSE estimator). Then
$\mathrm{MSE}_k=(1-b_k)^2+\Var\nu_k=1-b_k$, which gives $\Var\nu_k=b_k(1-b_k)$.
\end{proof}
The block outputs the unbiased symbols and writes $(1-b_k)/b_k$ to \meta{noise_var}, so the
soft demapper receives correctly scaled LLRs.
\codeexcerpt{mimo_linear}{ZF and unbiased MMSE detection with post-detection noise variances.}

\subsection{Maximum likelihood}
$\hat{\vect x}=\argmin_{\vect x\in\mathcal S^{N_t}}\|\vect y-\mat H\vect x\|^2$ is optimal and
achieves diversity $N_r$ per stream. The block searches all $M^{N_t}$ candidates exhaustively
(limited to 4096) and outputs hard decisions.

\begin{figure}[htbp]
\centering
\begin{tikzpicture}
\begin{semilogyaxis}[berplot,width=10cm,height=6.5cm,xmin=0,xmax=20,ymin=1e-5,ymax=0.5]
  \addplot[simline=accent] table[x=ebn0,y=zf]{\datafile{mimo_smux}};
  \addplot[simline={accent3,mark=square*}] table[x=ebn0,y=mmse]{\datafile{mimo_smux}};
  \addplot[simline={accent2,mark=triangle*}] table[x=ebn0,y=ml]{\datafile{mimo_smux}};
  \legend{ZF,MMSE,ML}
\end{semilogyaxis}
\end{tikzpicture}
\caption{$2\times2$ spatial multiplexing of QPSK over fast Rayleigh fading (simulated).
ZF and MMSE have diversity 1, the same slope as single-antenna Rayleigh fading, and MMSE gains
about 2\,dB. ML reaches diversity 2.}
\label{fig:smux}
\end{figure}

% ------------------------------------------------------------------
\section{Diversity}\label{sec:diversity}
\subsection{Receive diversity: maximal-ratio combining}
With one transmit antenna, ZF and MMSE reduce to maximal-ratio combining,
$\hat x=\vect h\Her\vect y/\|\vect h\|^2$, whose SNR is the sum of the branch SNRs. For BPSK with
$L$ i.i.d. Rayleigh branches of mean SNR $\bar\gamma$ each~\cite{proakis2008},
\begin{equation}
  \bar P_b=p^L\sum_{l=0}^{L-1}\binom{L-1+l}{l}(1-p)^l,\qquad
  p=\frac{1-\mu}{2},\quad \mu=\sqrt{\frac{\bar\gamma}{1+\bar\gamma}},
  \label{eq:mrc}
\end{equation}
and asymptotically $\bar P_b\propto\bar\gamma^{-L}$: the slope of the curve is the diversity order.

\subsection{Transmit diversity: the Alamouti code}
Two symbols are sent over two antennas in two symbol periods~\cite{alamouti1998}:
\begin{equation}
  \begin{array}{c|cc}
    & \text{antenna 1} & \text{antenna 2}\\ \hline
    t_1 & s_1 & s_2\\
    t_2 & -s_2\conj & s_1\conj
  \end{array}
  \qquad\Longrightarrow\qquad
  \begin{aligned}
    \tilde s_1&=h_1\conj r_1+h_2r_2\conj=(|h_1|^2+|h_2|^2)\,s_1+\tilde w_1,\\
    \tilde s_2&=h_2\conj r_1-h_1r_2\conj=(|h_1|^2+|h_2|^2)\,s_2+\tilde w_2 .
  \end{aligned}
\end{equation}
The code is orthogonal, so the linear combiner decouples the symbols and gives full diversity
2 without channel knowledge at the transmitter. The decoder sums over any number of receive
antennas. Because the transmit power is split between the antennas, the SNR is $3$\,dB below
two-branch MRC: $\bar P_b^{\text{Alamouti}}(\bar\gamma)=\bar P_b^{\text{MRC},L=2}(\bar\gamma/2)$.
\codeexcerpt{alamouti_dec}{Alamouti combining over $N_r$ receive antennas.}

\begin{figure}[htbp]
\centering
\begin{tikzpicture}
\begin{semilogyaxis}[berplot,width=11cm,height=7cm,xmin=0,xmax=22,ymin=1e-5,ymax=0.5,legend columns=2]
  \addplot[th=black!50,dashed] table[x=ebn0,y=awgn]{\datafile{diversity_theory}};
  \addlegendentry{AWGN}
  \addplot[th=accent] table[x=ebn0,y=siso]{\datafile{diversity_theory}};
  \addlegendentry{Rayleigh, theory}
  \addplot[sim=accent] table[x=ebn0,y=siso]{\datafile{diversity_sim}};
  \addlegendentry{SISO simulated}
  \addplot[th=accent3] table[x=ebn0,y=mrc2]{\datafile{diversity_theory}};
  \addlegendentry{MRC $L{=}2$, theory}
  \addplot[sim={accent3,mark=square*}] table[x=ebn0,y=mrc2]{\datafile{diversity_sim}};
  \addlegendentry{$1{\times}2$ simulated}
  \addplot[th=accent2] table[x=ebn0,y=alamouti]{\datafile{diversity_theory}};
  \addlegendentry{Alamouti, theory}
  \addplot[sim={accent2,mark=triangle*}] table[x=ebn0,y=alamouti]{\datafile{diversity_sim}};
  \addlegendentry{$2{\times}1$ simulated}
\end{semilogyaxis}
\end{tikzpicture}
\caption{BPSK over Rayleigh fading: single antenna (\cref{eq:rayleigh-ber}), $1\times2$ MRC and
$2\times1$ Alamouti (\cref{eq:mrc}). The $E_b/N_0$ axis is per receive antenna. Diversity 2
doubles the slope; Alamouti runs parallel to MRC, 3\,dB to the right.}
\label{fig:diversity}
\end{figure}
